\section{Planner Input and Action-Space Diagnostics}
\label{app:ablation}

\paragraph{Question and experimental unit.}
This compact ablation study documents behavior under changes to recent event
context and open-proposal access in the centralized execution harness described
in main-paper Section~\ref{main-sec:extended}. Its counters characterize the
implemented planner and handlers; the main paper's mechanism claim concerns
how generated state constrains actions and retains their effects. Aggregate
execution counts are not a measure of the mechanism's novelty or overall value. It uses the existing Clockwork Rain
Conservatory and Tidal Embassy worlds. Each world has one frozen initial agent
state and one fixed set of rules, reused across three configurations and three
stochastic repeats. Gemini 2.5 Flash is the only planning model. The resulting
18 trajectories are distinct from the earlier cross-model audit and its
model-specific generated worlds.

\paragraph{Conditions.}
Full supplies the current state, generated action catalog, current tasks,
current scripted intervention, and the 24 most recent logged events.
No event history sets only that last event list to empty. Current locations,
inventories, status summaries, tasks, and the current intervention remain
available. This is an input-history ablation, not removal of all historical
information or a test of decentralized episodic memory. Catalog only keeps
state, tasks, and history but removes the proposal option from the prompt and
JSON schema; the execution boundary also rejects proposals. Catalog actions
and adjacent-room movement remain available. As implemented in this harness,
catalog actions update statuses and relationships; they do not transfer or
create inventory. Its inventory limitation is therefore partly structural.

The earlier prospective Tidal configuration is not used as the manipulation:
its zero-valued memory limits are clamped to at least one by the legacy helper,
and it also changes task/cohort limits. Directly masking the planner's event
list avoids treating those settings as a clean memory ablation. Tests verify
that all other supplied fields remain identical when constructing full and
history-masked prompts from the same state.

\paragraph{Common computation and scheduling.}
Every trajectory has 24 rounds and six planning calls, at rounds 1, 5, 9, 13,
17, and 21. Each call requests 12 actions, three for each of four rounds.
An AI character may act at most once in each round; the scripted human is a
target. Calls use temperature 0.45, disabled thinking, a 6,144-token output
cap, and a 90-second timeout. At most two API attempts, including retries for invalid JSON, are allowed.
There is no alternate-model fallback. All conditions have the same call and
output limits; removing input history reduces input tokens, so total token
use is not equalized. Repeats are new stochastic API calls, not claims of a
provider-controlled random seed.

The schedule includes the existing human request, rejected unowned transfer,
resource shock, cross-room request, and human follow-up. Intervention rounds
and rules are shared; their recipient can depend on the state reached within
each trajectory. Current request records remain visible in every condition.
No actual participant is interacting during this evaluation. A fixed local
randomization orders the 18 run jobs, with up to four trajectories processed
concurrently. There is no adaptive stopping based on condition performance.

\paragraph{Protocol and retained outputs.}
The implementation was checked using three separate two-block pilot runs on
Clockwork. These six pilot calls are excluded. The study protocol, all input
states, and the runtime source snapshot were then hashed before the 108 formal
planning calls. This is a local pre-execution specification, not an external
preregistration. Model identifiers describe the requested API configuration;
an immutable provider snapshot is not claimed.

Every block preserves the supplied prompt payload and hash, JSON schema,
returned plan, call telemetry, executed events, and resulting state. Failures
are retained. The source snapshot contains the actual coordinator and action
handlers used by the experiment. The replay verifier reconstructs interventions
and re-executes every recorded action, checks the resulting state and event
hashes, and independently recomputes the reported counts without further LLM
calls.

\paragraph{Metrics and interpretation.}
The primary denominators are the 72 scheduled action slots per trajectory,
including slots lost to a planning failure. We report successful executions
and successful actions that change inventory or location. The latter requires
an actual difference in the sorted per-agent inventory/location snapshot;
relationship changes, status updates, and descriptions of intended effects
alone do not count. Inventory and movement counts are shown separately.
A location change does not by itself establish useful task progress. We also
retain proposal counts, complete-block counts, failed-state changes, and
provider-reported token usage. No weighted overall-quality score is formed.

Comparisons are paired within world, model, and repeat index. Repeats share
the same two worlds, and actions within a trajectory are dependent. We report
raw counts and descriptive differences, without treating the 1,296 action
slots as independent samples or claiming a population effect. The study does
not establish decentralized social coordination, human enjoyment, or a
universal benefit from longer history.

\begin{table}[t]
\centering\small
\setlength{\tabcolsep}{4pt}
\begin{tabular}{@{}lrrr@{}}
\toprule
Configuration & Success & Material & Inventory \\
\midrule
Full & 349 & 196 & 14 \\
No event history & 352 & 177 & 24 \\
Catalog only & 375 & 193 & 0 \\
\bottomrule
\end{tabular}

\caption{Aggregate planner diagnostics on two fixed worlds. Each row has
432 scheduled action slots across six trajectories. Success counts executed
actions; Material counts actual inventory or location changes among successes;
Inventory is its inventory-changing subset. Status and relationship changes
alone do not count as material.}
\label{tab:ablations}
\end{table}

\begin{table*}[t]
\centering\small
\setlength{\tabcolsep}{4pt}
\begin{tabular*}{\textwidth}{@{\extracolsep{\fill}}llrrrrrr@{}}
\toprule
World & Configuration & Rep. & Success & Material & Inventory & Move & Proposals \\
\midrule
Clockwork & Full & 1 & 64/72 & 44 & 1 & 43 & 1 \\
Clockwork & Full & 2 & 57/72 & 32 & 4 & 28 & 4 \\
Clockwork & Full & 3 & 55/72 & 31 & 1 & 30 & 3 \\
Clockwork & No event history & 1 & 59/72 & 30 & 5 & 25 & 11 \\
Clockwork & No event history & 2 & 55/72 & 38 & 8 & 30 & 9 \\
Clockwork & No event history & 3 & 63/72 & 29 & 8 & 21 & 22 \\
Clockwork & Catalog only & 1 & 68/72 & 43 & 0 & 43 & 0 \\
Clockwork & Catalog only & 2 & 68/72 & 27 & 0 & 27 & 0 \\
Clockwork & Catalog only & 3 & 60/72 & 30 & 0 & 30 & 0 \\
\midrule
Tidal & Full & 1 & 67/72 & 37 & 6 & 31 & 6 \\
Tidal & Full & 2 & 58/72 & 24 & 2 & 22 & 2 \\
Tidal & Full & 3 & 48/72 & 28 & 0 & 28 & 14 \\
Tidal & No event history & 1 & 56/72 & 20 & 3 & 17 & 26 \\
Tidal & No event history & 2 & 58/72 & 32 & 0 & 32 & 19 \\
Tidal & No event history & 3 & 61/72 & 28 & 0 & 28 & 6 \\
Tidal & Catalog only & 1 & 65/72 & 31 & 0 & 31 & 0 \\
Tidal & Catalog only & 2 & 53/72 & 33 & 0 & 33 & 0 \\
Tidal & Catalog only & 3 & 61/72 & 29 & 0 & 29 & 0 \\
\bottomrule
\end{tabular*}

\caption{All 18 formal trajectories. Each trajectory has 72 planned action
slots. Material counts successful actions with actual inventory or location
changes; Inventory and Move separate these consequences. Proposals counts
successfully executed proposals. Repeats share a frozen world and are not
independent world samples. Pilot runs are excluded.}
\label{tab:ablation-detail}
\end{table*}

\paragraph{Observed execution and cost.}
Table~\ref{tab:ablations} retains aggregate counts for every condition.
Full, No event history, and Catalog only execute 349, 352, and 375 of their
432 scheduled actions, respectively. Their inventory/location-changing
successes number 196, 177, and 193; inventory-changing subsets number 14,
24, and zero. The catalog-only material changes are all moves, consistent
with its restricted handlers. These totals do not establish a benefit from
supplied event history or measure successful completion of the scripted tasks.

All 108 formal planning blocks returned valid 12-action plans. One truncated
response required a second attempt, giving 109 API attempts. The 1,296 planned
slots produced 1,076 successful actions and 220 failures; every failure leaves
the recorded state unchanged. Of 566 inventory/location-changing successes,
528 are moves and 38 change inventory. Full executes 30 of 57 proposals;
No event history executes 93 of 138. These counts do not measure task utility.

Provider telemetry records 1,821,701 input tokens and 201,841 output tokens
across formal attempts, including the truncated attempt. Mean total tokens per
trajectory are 137,731 for Full, 74,598 for No event history, and 124,928 for
Catalog only. Cache-token counts are retained separately as a subset of input
usage, not added again. The six technical-pilot calls are outside these totals.

Relative to Full, the mean paired execution difference is $+0.7$ percentage
points for No event history, with repeat differences from $-15.3$ to $+18.1$;
its material-change difference is $-4.4$ points, ranging from $-23.6$ to $+11.1$.
Catalog only has a mean execution difference of $+6.0$ points (range $-6.9$ to
$+18.1$) and a material-change difference of $-0.7$ points (range $-8.3$ to
$+12.5$). These are descriptive ranges over six paired repeats on two worlds,
not confidence intervals. The variation and handler constraints limit the interpretation to these
recorded trajectories and effect types.

\paragraph{Trace provenance for the main-paper example.}
The example in main-paper Section~\ref{main-sec:extended} comes from the
Clockwork Full trajectory, repeat 1, event \texttt{model\_r09\_slot09}
(actual round 11). It is the first inventory-changing success in the first
lexicographically ordered Full trajectory with such an event, selected after
collection to illustrate the mechanism. In \texttt{block\_09.json}, Thistle
holds one Corrosive Inkwell and Elara holds none. The event records an approved
transfer of that item; \texttt{block\_13.json} reverses those holdings in the
supplied planner observation. The target-response field is unknown: this
outbound gift illustrates transfer checks, not independently observed
recipient acceptance. The same trajectory's scripted event
\texttt{coordinator\_rejection\_r09} rejects a transfer of an unowned item
without changing state. These records are included under
\texttt{llm\_ablation/runs/}; the illustrative selection adds no model calls
and does not replace the complete counts above.
